% s-080-p080-m002; original PDF 80, 81
\Aphorism{s-080-p080-m002}{73}
Im Gefolge einer Entscheidung, die aus vernünftiger Wahl hervorging, pflegt sich auch die Erkenntniss ihrer Bedingtheit, das Wissen um die Grenzen ihres Gültigkeitsbereiches einzufinden. Umgekehrt ist man den Sprüchen und Würfen des Zufalls gegenüber leicht geneigt, kritiklos zu werden und an ihre Allgemeinverbindlichkeit zu glauben — vielleicht mit der Schlussweise: da keine Vernunft darin ist, muss wohl etwas Höheres darin sein. So nimmt man sich die Freiheit, Berufe zu wechseln, aber beileibe nicht Religionen, warum? weil man zu seiner Religion kommt wie die Magd zum Kinde: zufällig, höchst zufällig! Derselbe Mensch, der im Übrigen sein Verhalten dem Einzelfall anpasst und mit der Relativität aller Interessen und Standpunkte rechnet, wird plötzlich fromm, feierlich, bornirt, wenn es sich um Vaterland oder politische Partei oder Standesvorurtheile handelt: warum? weil er sich von diesen Dingen, von seiner Zugehörigkeit zu ihnen keine Rechenschaft geben kann. Wir sind wahre Mystiker der Dummheit; je blinder, ungefährer, grundloser wir an eine Sache gerathen sind, desto hartnäckiger halten wir sie fest.
\par

